\chapter{\ifenglish Conclusions and Discussions\else บทสรุปและข้อเสนอแนะ\fi}

\section{\ifenglish Conclusions\else สรุปผล\fi}

จากการดำเนินการพัฒนาโมเดลสำหรับการพยากรณ์ปริมาณการผลิตไฟฟ้าจากแผงโซลาร์เซลล์โดยใช้ข้อมูลสภาพอากาศและข้อมูลการผลิตไฟฟ้าย้อนหลังจากแผงโซลาร์เซลล์เพื่อนำมาใช้เป็นข้อมูลในการฝึกฝนโมเดล 
โดยได้เลือกใช้เทคนิค Recurrent Neural Network (RNN) ประเภท Long Short-Term Memory (LSTM) ซึ่งเป็นเทคนิคที่เหมาะสมสำหรับการทำงานกับข้อมูลที่มีลักษณะเป็นลำดับเวลา (Time Series Data) 
หลังจากทำโมเดลและทดสอบพบว่าโมเดลสามารถพยากรณ์การผลิตพลังงานไฟฟ้าได้ในระดับที่น่าพอใจ โดยมีค่าประสิทธิภาพของโมเดลจากตัวชี้วัดทางสถิติ ได้แก่ ค่า $R^2$ เท่ากับ 0.8619 และมีค่า MAE 2.427 kWh 
ซึ่งแสดงให้เห็นว่าโมเดลมีความแม่นยำในการทำนายปริมาณการผลิตไฟฟ้าจากแผงโซลาร์เซลล์ได้ดี จากการทดสอบการใช้งานระบบพบว่าสามารถพยากรณ์ปริมาณการผลิตไฟฟ้าจากแผงโซลาร์เซลล์ได้ โดยมีค่าความคลาดเคลื่อนเฉลี่ย 4.574 kWhต่อจุดข้อมูล 
ซึ่งเป็นผลมาจากความแปรปรวนของสภาพอากาศ และข้อจำกัดของชุดข้อมูลที่ใช้ในการฝึกฝนโมเดล เนื่องจากข้อมูลสภาพอากาศที่นำมาฝึกฝนโมเดลไม่ได้ใช้ค่าที่มาจากสถานีวัดสภาพอากาศทีทำขึ้นแต่ตอนทดสอบใช้สภาพอากาศจากสถานีในการทำการทดสอบ
จึงทำให้ค่าที่ทำนายการผลิตไฟฟ้านั้นเกิดความคลาดเคลื่อน ทั้งนี้ ระบบดังกล่าวสามารถนำไปประยุกต์ใช้ในการวางแผนการบริหารจัดการพลังงานภายในอาคารได้อย่างมีประสิทธิภาพ 

\section{\ifenglish Challenges\else ปัญหาที่พบและแนวทางการแก้ไข\fi}

ในการทำโครงงานนี้ พบว่าเกิดปัญหาหลักๆ ดังนี้

\begin{enumerate}
        \item ข้อมูลค่าการผลิตไฟฟ้าย้อนหลังจากโซลาร์เซลล์บางช่วงมีค่าขาดหายไป หรือมีความไม่ต่อเนื่องของข้อมูล ทำให้ต้องมีการทำ Data Cleaning และจัดรูปแบบข้อมูลใหม่ก่อนนำไปใช้ในการฝึกฝนโมเดล
        \item รูปแบบของข้อมูลมีความแตกต่างกันเช่น ข้อมูลสภาพอากาศและข้อมูลการผลิตไฟฟ้าจากโซลาร์เซลล์ มีรูปแบบ Timestamp ที่แตกต่างกันทำให้ต้องมีการปรับรูปแบบของข้อมูลให้มีรูปแบบ Timestamp ที่เหมือนกันทั้งข้อมูลสภาพอากาศและข้อมูลการผลิตไฟฟ้า และทำการรวมข้อมูลเพื่อให้สามารถนำมาใช้ร่วมกันในโมเดลได้
        \item เนื่องจากข้อมูลการผลิตไฟฟ้าย้อนหลังจากโซลาร์เซลล์ที่มีเป็นแบบรายวันแล้วเราต้องการที่จะทำนายปริมาณการผลิตไฟฟ้าแบบรายชั่วโมง จึงต้องมีการทำการแปลงข้อมูลโดยการใช้เทคนิคการแปลงข้อมูลจากข้อมูลรายวันมาเป็นค่าการผลิตไฟฟ้าในแต่ละชั่วโมง
        \item เนื่องจากพึ่งติดตั้งสถานีวัดสภาพอากาศขนาดย่อมทำให้มีข้อมูลสภาพอากาศย้อนหลังไม่มากนัก ทำให้ต้องมีการใช้ข้อมูลสภาพอากาศจากแหล่งอื่นเพื่อใช้ในการโมเดล
\end{enumerate}   


\section{\ifenglish%
Suggestions and further improvements
\else%
ข้อเสนอแนะและแนวทางการพัฒนาต่อ
\fi
}

ข้อเสนอแนะเพื่อพัฒนาโครงงานนี้ต่อไป มีดังนี้

\begin{enumerate}
        \item เพิ่มปริมาณข้อมูลสำหรับการฝึกสอนโมเดลให้มากขึ้นเพื่อให้โมเดลมีความแม่นยำมากขึ้นในการทำนายปริมาณการผลิตไฟฟ้า
        \item พัฒนาให้ระบบสามารถทำ Real-time Prediction ได้จริง
        \item สามารถนำโมเดลที่พัฒนาขึ้นไปประยุกต์ใช้ร่วมกับระบบบริหารจัดการพลังงานเพื่อช่วยวางแผนการใช้พลังงานไฟฟ้าให้เกิดประสิทธิภาพสูงสุด
        \item เพิ่มระบบแจ้งเตือนเมื่อมีเซ็นเซอร์ตัวใดตัวหนึ่งมีปัญหาในการทำงานเพื่อให้สามารถแก้ไขได้อย่างรวดเร็ว
\end{enumerate}  